\documentclass[12pt,a4paper]{article}
\usepackage[utf8]{inputenc}
\usepackage[T1]{fontenc}
\usepackage{amsmath,amssymb,amsthm}
\usepackage{physics}
\usepackage{geometry}
\usepackage{enumitem}
\usepackage{hyperref}
\usepackage{fancyhdr}
\usepackage{titlesec}

\geometry{margin=1in}
\pagestyle{fancy}
\fancyhf{}
\rhead{Lecture 6}
\lhead{Advanced Quantum Mechanics}
\cfoot{\thepage}

\titleformat{\section}{\large\bfseries}{}{0em}{}
\titleformat{\subsection}{\normalsize\bfseries}{}{0em}{}

\newtheorem{definition}{Definition}

\title{\textbf{Lecture 6: Second Quantization} \\ \large From Harmonic Oscillators to Quantum Fields}
\author{Based on Leonard Susskind's \textit{Advanced Quantum Mechanics}}
\date{}

\begin{document}
\maketitle
\thispagestyle{fancy}

% ============================================================
\section{1. Review of the Harmonic Oscillator Algebra}
% ============================================================

The algebraic structure of a single harmonic oscillator:
\[
    [a, a^\dagger] = 1, \qquad N = a^\dagger a, \qquad H = \hbar\omega\left(N + \tfrac{1}{2}\right).
\]
The states $\ket{0}, \ket{1}, \ket{2}, \ldots$ are the occupation-number eigenstates, with $a^\dagger\ket{n} = \sqrt{n+1}\,\ket{n+1}$ and $a\ket{n} = \sqrt{n}\,\ket{n-1}$.

% ============================================================
\section{2. Many Independent Oscillators}
% ============================================================

Now consider a collection of independent harmonic oscillators, each labeled by an index $i$ and having its own frequency $\omega_i$.  Each oscillator has its own creation and annihilation operators:
\[
    [a_i, a_j^\dagger] = \delta_{ij}, \qquad [a_i, a_j] = 0, \qquad [a_i^\dagger, a_j^\dagger] = 0.
\]
The total Hamiltonian is:
\[
    H = \sum_i \hbar\omega_i\left(a_i^\dagger a_i + \frac{1}{2}\right).
\]
A general state is specified by the occupation numbers of each oscillator:
\[
    \ket{n_1, n_2, n_3, \ldots} = \frac{(a_1^\dagger)^{n_1}}{\sqrt{n_1!}}\frac{(a_2^\dagger)^{n_2}}{\sqrt{n_2!}}\cdots\ket{0, 0, 0, \ldots}.
\]
This is the \textbf{Fock space} representation.

% ============================================================
\section{3. Reinterpretation: Particles as Excitations}
% ============================================================

The key insight of second quantization is a \emph{change of language}:
\begin{center}
\begin{tabular}{l l}
\textbf{Oscillator language} & \textbf{Particle language} \\
\hline
Oscillator $i$ in $n$-th excited state & $n$ particles in single-particle state $i$ \\
$a_i^\dagger$ raises excitation of oscillator $i$ & $a_i^\dagger$ creates a particle in state $i$ \\
$a_i$ lowers excitation of oscillator $i$ & $a_i$ annihilates a particle in state $i$ \\
$N_i = a_i^\dagger a_i$ counts excitation level & $N_i$ counts particles in state $i$ \\
Vacuum $\ket{0}$ (all oscillators in ground state) & Empty space (no particles) \\
\end{tabular}
\end{center}

\medskip
The mathematics is identical; only the interpretation changes.  The occupation number $n_i$ is now the number of particles in single-particle state $i$.

\subsection{3.1 Total particle number}

The total number of particles is:
\[
    \hat{N}_{\text{total}} = \sum_i a_i^\dagger a_i.
\]
This operator counts the sum of all occupation numbers.

% ============================================================
\section{4. Bosonic Many-Particle States}
% ============================================================

For bosons, the commutation relations $[a_i, a_j^\dagger] = \delta_{ij}$ allow any number of particles in each state.  The symmetry of the wave function under particle exchange is automatically guaranteed by the commutation of creation operators:
\[
    a_i^\dagger a_j^\dagger = a_j^\dagger a_i^\dagger.
\]
Creating particle $i$ then particle $j$ is the same as creating $j$ then $i$---the state is automatically symmetric.

% ============================================================
\section{5. Creation and Annihilation Operators for Particles}
% ============================================================

\begin{definition}[Creation operator]
The operator $a_i^\dagger$ \textbf{creates} a particle in single-particle state $i$:
\[
    a_i^\dagger\ket{n_1, \ldots, n_i, \ldots} = \sqrt{n_i + 1}\,\ket{n_1, \ldots, n_i + 1, \ldots}.
\]
\end{definition}

\begin{definition}[Annihilation operator]
The operator $a_i$ \textbf{annihilates} (removes) a particle from single-particle state $i$:
\[
    a_i\ket{n_1, \ldots, n_i, \ldots} = \sqrt{n_i}\,\ket{n_1, \ldots, n_i - 1, \ldots}.
\]
If there is no particle in state $i$ ($n_i = 0$), then $a_i$ gives zero.
\end{definition}

% ============================================================
\section{6. Why ``Second Quantization''?}
% ============================================================

The term ``second quantization'' is somewhat misleading.  There is really only one quantization---the standard rules of quantum mechanics.  The name arose historically because the procedure \emph{looks like} quantizing the already-quantized single-particle wave function by promoting it to an operator.

\medskip
The real power of the formalism is:
\begin{enumerate}[nosep]
    \item It automatically handles identical-particle symmetry (symmetric for bosons, antisymmetric for fermions).
    \item It naturally describes systems with a \emph{variable} number of particles (creation and annihilation).
    \item It provides the natural language for quantum field theory.
\end{enumerate}

\medskip
\noindent\textit{In the next lecture, we will construct the quantum field operator $\Psi(x)$ as a superposition of annihilation operators weighted by single-particle wave functions, completing the bridge from particles to fields.}

\end{document}
